\chapter{Evaluation setup and research material}
\label{ch:evaluation}

This chapter describes how the analysis of Chapter~\ref{ch:analysis} is evaluated: the corpus
(Section~\ref{sec:corpus}), every ratio and measure (Section~\ref{sec:metrics}), and the five
evaluation methods with what each can and cannot show (Section~\ref{sec:protocol}).
Chapter~\ref{ch:results} reports the results.

\section{Corpus selection and validity}
\label{sec:corpus}

\subsection{Selection}

Ten candidate applications were screened against six criteria: a real application rather than a
library or template; a nested, lazily loaded route tree; guards present; end-to-end test
configuration for a browser driver; retrievable history in which tests and source change together;
and an active repository. Six candidates were retained (Table~\ref{tab:corpus}), each fixed to a
specific commit so that the measurements can be repeated on the same code.

Criterion (4) was checked as configuration, not as a usable suite. At the pinned commits only
alfresco has specifications the test-source cross-check can read, so that method covers alfresco
alone. Ghostfolio and bitwarden carry no browser-test configuration at all, so for them the
criterion was not met. It was therefore not applied uniformly: \texttt{gravitee} was rejected
for a thin UI suite while two applications with none were kept.

\begin{table}[t]\small
\caption{The corpus: commit, route count and role in the study of each application.}
\label{tab:corpus}
\begin{tabularx}{\textwidth}{@{}lYrY@{}}
\toprule
Application & Commit & Routes & Role in the study \\
\midrule
gothinkster/angular-realworld-example-app & \texttt{dd99ed2c} & \RoundTwoRoutesRealworld{} & size floor; all literal paths \\
alfresco/alfresco-content-app & \texttt{b88d47f4} & \RoundTwoRoutesAlfresco{} & 87 Playwright specs including a navigation suite \\
ghostfolio/ghostfolio & \texttt{9884fa0b} & \RoundTwoRoutesGhostfolio{} & constant-valued route paths; self-hostable \\
kubernetes/dashboard & \texttt{0ba796d9} & \RoundTwoRoutesKdash{} & mid-size control \\
bitwarden/clients (web) & \texttt{e40ce6f0} & \RoundTwoRoutesBitwarden{} & many guards, guard factories among them \\
openMF/web-app & \texttt{7f781796} & \RoundTwoRoutesMifos{} & size ceiling \\
\bottomrule
\end{tabularx}
\par\smallskip{\footnotesize\noindent\emph{Note.} Routes are the route entries reached by the independent route parser of Section~\ref{sec:protocol}, as it stands in the final version (Section~\ref{sec:res-round2}). A count marked $\geq$ is a lower bound: the parser records a blind spot for that application, so the entries behind it are not in the figure; Table~\ref{tab:coverage} gives how many blind spots each application has. The marks are on \RoundTwoLowerBoundWhere{}, on these grounds: \RoundTwoLowerBoundReasons{}.}
\end{table}

Four candidates were rejected; the counts were read from the repositories and are not generated.
\texttt{angular/components} is a component library with no product router. \texttt{vendure} has 169
end-to-end specifications, all of which exercise the GraphQL server and none the user interface.
\texttt{gravitee} declares 618 routes but has 14 UI specifications. \texttt{SAP/spartacus} has a
large Cypress suite, but its routing is configuration-driven in a way the extractor does not
support; that reason is tool support, which is not one of the six criteria, so this rejection rests
on a ground added after screening began.

Route-tree size and end-to-end test density were inversely related: no application provides both a
large route tree and execution-based ground truth. The methods that need test suites are applied
only to the applications that have them, and the other methods to all six.

\subsection{Validity of an open-source corpus}

Kalliamvakou et al.\ show that many GitHub repositories are inactive, personal or not software
projects, and recommend filtering candidates by explicit criteria and reporting them
\cite{kalliamvakou2014promises,kalliamvakou2016indepth}, which the selection above does. The
remaining threats are external validity (open-source applications may differ from proprietary ones)
and a possible selection bias, since I selected the corpus and wrote the extractor. The selection
protocol is therefore reported in full: the criteria were fixed before any results were obtained,
but one candidate was rejected on tool support and two retained applications do not meet the
test-suite criterion. The corpus spans nearly two orders of magnitude in routes, and the commits
and harness are released so the measurement can be repeated.

None of this addresses a further threat. The extractor was developed on this corpus and is measured on it
(Section~\ref{sec:limitations}), so the figures are fitted to it. The one measurement outside it runs the frozen extractor on the seven
subjects of Rima et al.\ against their ground truth, preregistered before any subject was mapped (Section~\ref{sec:res-rima}). A case study on a proprietary
Angular application, reporting only aggregate figures, is planned as the most direct response
(Chapter~\ref{ch:discussion}). The literature reviewed in Chapter~\ref{ch:background} offers no
standard benchmark corpus of single-page applications either, so the corpus and its manifest are
an output of the thesis.

\section{Metrics}
\label{sec:metrics}

Each research question is answered with one main measure. Table~\ref{tab:ratios} lists every ratio
the thesis reports, and the paragraphs after it define the terms the table uses.

\begin{table}[t]\small
\caption{The ratios and measures of the thesis, with numerator, denominator and research question.}
\label{tab:ratios}
\begin{tabularx}{\textwidth}{@{}>{\raggedright\arraybackslash}p{2.4cm}YY>{\raggedright\arraybackslash}p{0.9cm}>{\raggedright\arraybackslash}p{2.0cm}@{}}
\toprule
Name & Numerator & Denominator & RQ & Reported in \\
\midrule
Node coverage & declared screens that have a node & screens the independent route parser declares & RQ1 & \ref{sec:res-baseline} \\
Must conversion & in-scope constructs witnessed by a must edge & in-scope constructs & RQ1 & \ref{sec:res-baseline} \\
Recovered share & in-scope constructs witnessed by an edge of any tier & in-scope constructs & RQ1, RQ3 & \ref{sec:res-after}, \ref{sec:res-round2} \\
Accounted share & recovered constructs plus in-scope constructs recorded as dynamic or unplaced & in-scope constructs & RQ2 & \ref{sec:res-defects} \\
Must-tier precision & must edges whose witness passes (lower bound) or does not fail (upper bound) & all must edges & RQ3 & \ref{sec:res-soundness} \\
Node-level recall & distinct URLs named in an application's end-to-end test sources that have a node & distinct URLs so named & RQ1 & \ref{sec:res-baseline} \\
Structural-reference recall and precision & (site, screen) pairs the reference and a level share & pairs in the reference (recall) or pairs the level places (precision) & RQ3 & \ref{sec:res-groundtruth} \\
End-to-end containment & observed screen pairs joined by an edge of the graph & observed screen pairs & RQ3 & \ref{sec:res-e2e} \\
Scope width & mean screens a site is attributed to & the same when every site has scope $\top$ & RQ3 & \ref{sec:res-width} \\
Detection rate & injected changes that produce the predicted change in the graph & injected changes the graph can score (excluded mutants reported) & RQ4 & \ref{sec:res-mutation} \\
Cause trace & defect classes traced to a cause in the source, by kind (read by the final version, still not read, beyond a static reader); no ratio & defect classes & RQ5 & \ref{sec:res-round2} \\
\bottomrule
\end{tabularx}
\par\smallskip{\footnotesize\noindent\emph{Note.} Each measure has one name, used throughout. A construct is \emph{in scope} when it is neither a query-only navigation nor excluded by the first rows of Table~\ref{tab:loss-rules}. Node coverage is the only measure called coverage in this thesis.}
\end{table}

\textbf{Navigation construct.} A navigation construct is a \texttt{routerLink} binding in a
template, or a call to \texttt{Router.navigate} or \texttt{Router.navigateByUrl} in class code,
in the application source excluding tests. Constructs are found by a separately written parser
that shares no code with the extractor. Table~\ref{tab:repairs} qualifies what that
independence covers: for the route-path forms added during development, the two resolve the same
constructs by the same rules, and three of the independent parser's corrections were made because it
disagreed with the extractor.

\textbf{Must conversion (RQ1).} Must conversion is the fraction of in-scope constructs that are
witnessed by at least one must edge. Three rules make it precise.
\begin{itemize}
\item A construct is matched to a witness on its own line or on a line either side, to absorb
offsets in the extractor's locations. Two constructs on adjacent lines can therefore both be
credited by one witness, so the ratio is an upper bound at line granularity.
\item Query-only navigations and the two exclusions of Table~\ref{tab:loss-rules} are left out of
the denominator as well as the numerator.
\item When several constructs produce the same edge, the extractor keeps the most certain as the
witness and records the others as merged witnesses, each with the modality it earns on its own. A
construct counts as converted only where it earns must itself.
\end{itemize}

\textbf{The other ratios.} More than one of these measures has been called coverage in the
literature this thesis builds on, but they are not interchangeable: on this corpus they differ by
more than an order of magnitude (\RoundTwoAccounted{} of sites accounted for against
\RoundTwoMustConv{} converted to must). Node coverage rises if less of an application is read, so
the blind-spot count is printed beside it. The recovered share widens the must-conversion numerator
to an edge of any tier. The accounted share alone measures whether the losses are declared, not how
much is recovered. Coverage without a qualifier means node coverage; the gap in recovered
constructs is called the recovery gap, and the loss tables measure it.

\textbf{Agreement between a scope and the rendered interface (RQ3).} The measures above cannot be
asked the question the contribution turns on: is the set of screens that attribution computes for a
site the set that site can really fire from? The \emph{rendered-interface check}, run on the harness
of the runtime route oracle (Section~\ref{sec:protocol}), answers one direction. The harness
navigates to each screen the graph holds, reads the in-app links the DOM carries and compares them
with the edges the graph's per-screen query yields there.

A link rendered at a screen where the graph has no such edge is unambiguous: the site was
attributed to too few screens or never read. An edge the graph has and the DOM does not show is
weaker, since a template may render links conditionally and this harness signs in to nothing; where
a target could not be resolved to one address, every address the static prefix admits is recorded at
reduced confidence, which is not a claim that the screen links there. This direction is therefore
judged only over predictions that name a destination exactly, and their number is published. Three
bounds apply. Only template-sourced sites are judged, since a \texttt{Router.navigate} call in an
event handler leaves nothing in the DOM. A screen whose guard redirects to a sign-in page without a
backend renders another screen's links and cannot be judged; such screens are counted apart, and on
this corpus they are most screens (Section~\ref{sec:res-scope}). Ghostfolio's per-locale build does
not fit the harness, so it is measured by every other method and by none of this one.

\textbf{The ablation and the structural reference (RQ3).} Two measurements ask whether the
attribution computes a set of screens that agrees more closely with a structural reference than a
simpler analysis does, or only a larger one.

The ablation withholds the mechanisms of Section~\ref{sec:attribution} one at a time and extracts
the whole corpus at each level:
\begin{itemize}
\item \emph{Owned}: a site belongs to the screen whose routed component's own file holds it, the
direct ownership of the Android literature.
\item \emph{Nesting}: adds component containment through template element selectors.
\item \emph{Rooted}: adds the shell rule, so a component the root renders is attributed to every
screen. Oshima and Kishi assemble their page information from the root component as well
(Section~\ref{sec:dynamic-spa}), so this level measures how far the prior work reaches here.
\item \emph{Outlet}: adds router-outlet descent.
\item \emph{Full}: adds the member-level call graph, with its fallback join and the stand-ins for
guards and effects.
\end{itemize}
Each level adds to the one before it, so the \emph{full} step is not one mechanism alone. A sixth
level, \emph{class}, runs beside the nested ones and replaces the call graph with the per-class rule,
which joins a navigating service over every component that injects its class. The harness checks on
the real graphs that the levels nest: a (site, screen) pair present at one level and absent at the
next is a defect in the switch, and no table is written.

The structural reference is scored in both directions, which no other method allows. For one
application, the screens each navigation site fires from were read off the sources by hand, with a
derivation per site, in \path{route-oracle/ground-truth/realworld.json}. It is the structural
reference of Section~\ref{sec:scope}, called the ground truth below. A pair the reference has and a
level does not is a scope too small by the reference; the converse is a scope too large. Four
limits apply.
\begin{itemize}
\item On realworld it is one application, small enough to read exhaustively. A second, alfresco, is read
under the same rules (Section~\ref{sec:res-groundtruth}).
\item It is the author's reading, so it is not independent of the tool's author. It is published with
its reasoning so a reader can check it.
\item Its rules are the analysis's own relations, and a template condition or guard that admits a
screen only in one sign-in state does not narrow a scope under them. The ground truth therefore
tests whether the implementation applies the relations correctly, not whether they match what the
running application renders.
\item It is checked in one direction only: the scorer reports every site the graph places that the
ground truth omits.
\end{itemize}

\textbf{The end-to-end navigation check (RQ3).} The structural reference applies the analysis's own
relations, so it cannot show that they miss a navigation the application performs. This check takes
its reference from execution. Two applications ship an end-to-end suite that can be run here.
The first, realworld, runs against the public demonstration backend, which is not pinned. The second, mifos, runs against a local
Fineract server with the image pinned by digest. A small module added to each application logs
every navigation the router completes while the suite runs, with the route patterns of the screen
left and the screen reached.

An observed pair is such a source and target pattern for a navigation the application itself
triggered; back-button navigations and the first screen after a full page load are excluded. A pair
is \emph{contained} if any edge, must or may, joins the two screens once scope nodes are expanded to
their members, with parameter names ignored. An edge to an unknown target never counts. A hit is any
such edge, not the edge of the site the test used, so the check does not test the scope of an
individual site. Redirected navigations are compared on the address they asked for, as
preregistered. The scorer's first departure from this rule is listed in
Table~\ref{tab:repairs}, and its figure is not used.

A missed pair is classified from the graph alone, by rules over the graph so that no assignment is
made by hand. The classes are tested in this order, and the first match wins:
\begin{itemize}
\item \emph{redirect}: the navigation was redirected and the graph holds the pair on the address it landed on (the preregistered redirect class);
\item \emph{base}: the target is reached from a screen above the source but not from the source, as a relative link resolved against the wrong route appears;
\item \emph{dynamic}: the source screen's routed component, or one above it, holds a site whose target the extractor records as unresolved;
\item \emph{other}: none of these.
\end{itemize} The base and dynamic classes were defined after the misses had been
read. Because of the order, the dynamic count is what is left after base. On mifos every miss met
the dynamic test whatever the order, so that count said little alone; Section~\ref{sec:res-e2e} names
the sites, which were repaired after the check.

The counting rules, three validity checks and predictions for each application were written down
before any log was read, except for that departure. The checks are that the instrumented and the
plain suite pass and fail the same tests, that at least ten distinct pairs are observed, and that a
second instrumented run observes the same pair set. The check measures recall on the transitions
the suite exercises and nothing else. The graph is an over-approximation, so that recall is
containment, and its weight depends on how many (screen, screen) pairs the graph holds, which is
reported beside it. A graph edge no test takes is not an error, so the check cannot measure
precision, and one run cannot refute a must edge.

\textbf{Loss breakdown (RQ2).} Each construct is matched to the edges whose witness cites it. A
construct that is in scope, cited by no must or may edge and not accounted for is \emph{lost}. It is
assigned to exactly one loss category by the first rule in Table~\ref{tab:loss-rules} that applies.
The categories therefore partition the lost constructs. Named outlets are removed before the
assignment, and guards are measured outside it. The categories Constant target, Node missing and
Other were not put through the admission rule that maps a category to a defect class
(Chapter~\ref{ch:results}).

Guard defects do not appear in the table, since a wrong guard mislabels an edge rather than removing
it. No method here checks a guard label directly; guard handling is exercised only by the
guard-removal mutation operator and by the source review of Section~\ref{sec:res-round2}. The guard
defect was found by red-team review of the graph and the source, not by any check above.

\begin{table}[t]\small
\caption{Rules that classify a construct, in the order the harness tests them.}
\label{tab:loss-rules}
\begin{tabularx}{\textwidth}{@{}lY@{}}
\toprule
Category & The construct \ldots \\
\midrule
Query only & changes only the query string of the current screen
(\texttt{navigate([], \{queryParams\})}), so it is not a transition and is not counted as
a loss \\
Named outlet & targets an auxiliary outlet (\texttt{\{outlets: \ldots\}}) \\
No commands & names no destination at all \\
Converted, must or may & is cited as witness by an edge of that tier \\
Accounted, dynamic & is cited by an edge into a sink: the extractor records its destination as unresolved \\
Accounted, unplaced & is recorded by the extractor as a site it read but could not place \\
Attribution, service & has no owning component, or its owner is a class that is not a
component: a service, guard, interceptor, resolver or effect. A site in a template whose
owner cannot be resolved falls to \emph{other} instead \\
Attribution, shared & is in a component that no route renders directly \\
\addlinespace
\multicolumn{2}{@{}l}{\emph{Owner is a route component, but that component has no node:}} \\
\quad Component unlinked & a node for the owner's route exists, but the extractor did not
link it to the component \\
\quad Constant path & every declaration of the owner's route gives its path as a constant
reference rather than a literal \\
\quad Other & some declaration of the owner's route does have a node \\
\quad Lazy boundary & some declaration of the owner's route is in a lazily loaded routes
file \\
\quad Node missing & a declaration exists but none of the above applies, or neither parser
found a route for the owner \\
\addlinespace
Constant target & owner has a node, but the construct's own target is not a literal \\
Other & matches none of the above; inspected by hand \\
\bottomrule
\end{tabularx}
\par\smallskip{\footnotesize\noindent\emph{Note.} The first three rows take a construct out of the denominator, the next counts it as recovered in the tier that cites it, and the two after that as accounted for. Only a construct that reaches the rows below is lost. The indented rules are the nested decision taken when the construct's owning route component has no node.}
\end{table}

The nested branch decides a category from the \emph{owning} route component, so a construct can be
assigned to a node category and still be converted once attribution gives it a scope through the
shell or a subtree (Section~\ref{sec:res-after}). The size of \emph{other} counts constructs left
for inspection by hand, not one cause.

\textbf{Change in recovery and must-tier precision (RQ3).} Must conversion and the loss breakdown are
measured on one commit of the extractor, run twice on identical inputs with the attribution pass
disabled and enabled. With the pass off the extractor reproduces the baseline edge set exactly on
every application (Section~\ref{sec:implementation}).

\emph{Must-tier precision} is the fraction of must edges that meet two conditions. The witness line contains a construct of
the claimed kind with the claimed target, \emph{and} the destination node sits at an address the
independent route set declares; otherwise a node minted at an address the application never mounts
would pass.

An edge whose destination the check cannot read is neither a pass nor a failure, so precision is
stated as bounds: the share of must edges that pass and the share that do not fail.
RQ3 is answered positively only if the recovered share and must conversion, over the same construct
denominator, do not fall, the recovered share rises, and must-tier precision does not fall. This is
read as no detected failure in either run, not as a precision that did not fall where the check
reads only part of the tier. The bounds of the two runs cover different sets of must edges, so
they are not compared as rates.

\textbf{Detection rate (RQ4).} The detection rate is the fraction of injected source changes
that produce the predicted change in the graph, reported per mutation operator with its $n$. A
change whose construct the graph does not contain is excluded and not counted as a miss, since it
can only repeat a recovery gap that the loss tables measure. The rate is therefore conditional on
recovery, and the number of excluded mutants is reported alongside it.

\textbf{Cause trace (RQ5).} Each defect class of Table~\ref{tab:defects} is traced to its cause in the
source, with one of three outcomes: the final version reads the cause (the loss was the tool's and
not the application's), it does not yet read it although a static reader could, or the cause lies
beyond a static reader of a route table. The result is a list of causes by outcome, not a ratio.
Every site that yields no transition still carries a recorded reason (Section~\ref{sec:res-round2}).

\section{Evaluation methods}
\label{sec:protocol}

Figure~\ref{fig:pipeline} shows where the methods act on the extraction pipeline. Five methods are used; RQ3 adds four further checks (the rendered-interface check, the ablation, the structural reference and the end-to-end check), which Table~\ref{tab:methods} does not list. None establishes correctness alone, so Table~\ref{tab:methods} states what
each can show and where it is blind. Figure~\ref{fig:oracles} shows which part of the graph each
method checks and where its ground truth comes from.

\begin{figure}[t]
\centering
\begin{tikzpicture}[font=\scriptsize, >=Stealth,
  stage/.style={draw=ink, rounded corners=2pt, fill=white, align=center, text width=22mm, minimum height=10mm, inner sep=2pt},
  opt/.style={stage, dashed, text width=27mm},
  check/.style={draw=ink, rounded corners=2pt, fill=mutedlight, align=left, inner sep=4pt, text width=42mm}]
  \node[stage] (a) at (0,5mm) {Angular\\source};
  \node[stage] (b) at (34mm,14mm) {route\\reader};
  \node[stage] (c) at (68mm,14mm) {screen\\nodes};
  \node[stage] (n) at (34mm,-4mm) {navigation\\sites};
  \node[stage, fill=accentlight] (e) at (68mm,-4mm) {edges: must, may,\\unknown};
  \node[opt] (s) at (68mm,-22mm) {scope attribution\\(optional pass)};
  \node[check] (m) at (112mm,10mm) {\textbf{Evaluation methods} (Table~\ref{tab:methods})\\[1pt]
    independent route parser $\cdot$ witness check $\cdot$ \mbox{test-source} cross-check $\cdot$ mutation testing $\cdot$ runtime route oracle};
  \node[check] (r) at (112mm,-20mm) {\textbf{RQ3 checks}\\[1pt] rendered-interface check $\cdot$ ablation $\cdot$ structural references $\cdot$ end-to-end suites};
  \draw[->] (a.east) -- ++(5mm,0) |- (b.west);
  \draw[->] (a.east) -- ++(5mm,0) |- (n.west);
  \draw[->] (b) -- (c);
  \draw[->] (n) -- (e);
  \draw[->, dashed] (s) -- (e);
  \draw[->, muted] (c.east) -- ([yshift=4mm]m.west);
  \draw[->, muted] (e.east) -- ++(7mm,0) |- ([yshift=-4mm]m.west);
  \draw[->, muted] (e.east) -- ++(7mm,0) |- ([yshift=4mm]r.west);
  \draw[->, muted] (s.east) -- ([yshift=-4mm]r.west -| s.east) -- ([yshift=-4mm]r.west);
\end{tikzpicture}
\caption{What is evaluated. The extractor reads route tables into screen nodes and navigation sites into edges, tiered must, may or unknown by the demotion rules; scope attribution, an optional pass, places sites that belong to no single screen. The five evaluation methods of Table~\ref{tab:methods} check nodes and edges against independent readings; the RQ3 checks examine the attribution.}
\label{fig:pipeline}
\end{figure}


Four of the five read the application's sources, and two readers of the same text can be wrong in
the same way. The fifth, the runtime route oracle, was added for that reason.

\begin{figure}[t]
\centering
\begin{tikzpicture}[node distance=9mm and 16mm]
  \node[box, hl, minimum width=30mm, minimum height=12mm] (g) {extracted\\graph};
  \node[box, above left=of g, minimum width=38mm, font=\footnotesize] (b) {\textbf{witness check}\\truth: the source line\\each edge cites};
  \node[box, above right=of g, minimum width=38mm, font=\footnotesize] (c) {\textbf{differential test}\\truth: an independent\\route parser};
  \node[box, below left=of g, minimum width=38mm, font=\footnotesize] (d) {\textbf{test-source cross-check}\\truth: URLs named in\\test sources};
  \node[box, below right=of g, minimum width=38mm, font=\footnotesize] (f) {\textbf{mutation}\\truth: the change\\that was injected};
  \draw[<->] (b) -- node[lbl, above right, pos=0.4] {must edges} (g);
  \draw[<->] (c) -- node[lbl, above left, pos=0.4] {nodes} (g);
  \draw[<->] (d) -- node[lbl, below right, pos=0.4] {nodes} (g);
  \draw[<->] (f) -- node[lbl, below left, pos=0.4] {graph delta} (g);
  \node[font=\scriptsize, text=muted, above=0.5mm of b] {blind to missing edges};
  \node[font=\scriptsize, text=muted, above=0.5mm of c] {independent, not stronger};
  \node[font=\scriptsize, text=muted, below=0.5mm of d] {named subset only, no test run};
  \node[font=\scriptsize, text=muted, below=0.5mm of f] {needs the construct in the graph};
\end{tikzpicture}
\caption{The four evaluation methods that read the applications' sources, the graph part each checks, and its ground truth.}
\label{fig:oracles}
\par\smallskip{\footnotesize\noindent\emph{Note.} The fifth method, the runtime route oracle, is not drawn because its ground truth is the running application rather than any part of the graph. No method sees everything; the grey notes give the blind spot of each, as listed in Table~\ref{tab:methods}.}
\end{figure}


\begingroup\small
% Taller than a page, so it breaks rather than overflowing one.
\begin{xltabular}{\textwidth}{@{}lYY@{}}
\caption{Evaluation methods, the quantity each measures, and the blind spot each has.}
\label{tab:methods}\\
\toprule
Method & Quantity & Blind spot \\
\midrule
\endfirsthead
\multicolumn{3}{@{}l}{\itshape Table~\ref{tab:methods} continued from the previous page}\\
\toprule
Method & Quantity & Blind spot \\
\midrule
\endhead
\midrule
\multicolumn{3}{r@{}}{\itshape continued on the next page}\\
\endfoot
\bottomrule
\endlastfoot
Witness check & Must-tier precision. For every must edge, the witness file is
re-parsed at the claimed line, the construct kind and target are compared, and the
destination node's address is tested against the independent route set. &
Cannot see missing edges. Target agreement is invalid for over-approximated may
edges, whose witness proves only that the call site exists. Where both tools read a table at the
wrong prefix, they agree on a wrong address (Section~\ref{sec:res-probe}). Target agreement is
literal, so an edge to the screen a redirect lands on is reported as a failure. \\
\addlinespace
Differential test & Disagreements with a separately written route-configuration
parser. &
Diagnostic, not a metric. The independent route parser is independent, not stronger, and for the path
forms added during development (Table~\ref{tab:repairs}) it is no longer independent of the extractor at
all. \\
\addlinespace
Test-source cross-check & Node-level recall against a fourth reader of the source. The URL literals in
the application's end-to-end specifications are collected; a URL its own test suite names is an
address its authors expect to exist, and the extractor is asked whether it minted a node there. &
Reads test source text and runs nothing, so a collected string is an address a test \emph{names}, not
one a test was observed to reach. Tested subset only. Textual extraction under-samples click-driven navigation. A direct \texttt{goto} is not a UI
transition, so visit pairs are not counted as edges. \\
\addlinespace
Mutation test & Detection rate over four operators: rename a path, remove a guard,
delete a \texttt{routerLink}, retarget a \texttt{navigate}. Ground truth is the
authored delta. Detection is judged on the node and edge sets and on the witness set,
since several constructs can produce one edge. &
A mutant whose construct is absent from the baseline graph re-detects a recovery gap and
is excluded, as is a guard mutant on a screen that no edge enters. A component several routes render maps to every screen those routes declare, which errs towards scoring a mutant rather than dropping it. Detection is conditional on
recovery, and the attributed and unattributed graphs do not exclude the same mutants, so they are
also compared over the mutants both can score, paired by operator, file, line and subject. The
sample is drawn from the graph, so it is not independent of the artefact under test. \\
\addlinespace
Runtime route oracle & Node-level soundness against the framework itself. Each application is built in
its development configuration, loaded in a headless browser, and asked for the route configuration the
Angular router holds, with every lazily loaded subtree expanded through the application's own
\texttt{loadChildren}. The addresses the extractor mints are compared with the addresses the framework
mounts. Nothing re-reads the route files. Which addresses count as mounted is itself a reading, so the
harness navigates the running application to a list of addresses and records what the recogniser made
of each, before any \texttt{canActivate} guard runs and after the \texttt{canLoad} and
\texttt{canMatch} guards that take part in matching. &
Requires the application to build and boot without its server. It reads the configuration and not the
rendered screen, so a route that a guard always refuses still counts as mounted. An address claimed
through a \texttt{UrlMatcher} cannot be decided from the configuration. A workspace application
that composes the measured library mounts the shell's own routes as well, so part of any difference
is scope and not reading. Its result files record no extractor commit or runtime version
(Section~\ref{sec:reproducibility}). \\
\end{xltabular}
\endgroup

\subsection{Preconditions}

Two properties of the extractor were checked before any measurement. \emph{Determinism}: two
extractions from the same unchanged source give identical nodes and edges. \emph{No accumulated
state}: extracting into an existing database gives the same graph as extracting after deleting it.
Without them, differences seen in mutation testing could not be told from noise.

\subsection{Corrections made to the protocol}
\label{sec:corrections}
\label{app:repairs}
\label{app:validation}

The final figures are those of the final version of the protocol. The extractor was developed within this thesis, and three versions are named: the \emph{early development version} (commit \EarlierCommit{}), before the corrections the evaluation prompted; the \emph{intermediate version} (extractor \RoundExtractor{}, with a must tier of \RoundMustEdges{} edges, tables as printed at \RoundCommit{}); and the \emph{final version}, which the chapters present. Of the intermediate version only the must-tier figures are cited, as printed; the runs behind the printed tables were lost, and an archived re-run of the same extractor gave a different loss ledger. The figures of superseded instruments cannot be regenerated with the shipped harness, and none is cited as a result. A surprising
figure is checked against its inputs before it is written down, since such a figure has more than
once proved a property of the measurement and not of the extractor.

The mutation harness was corrected during development after its misses had been read, so the final
run is not an independent test. The final run used the same seed as the run whose misses were read
but a partly different sample. \MutOverlap{} of its \MutSampled{} mutants are at the same location
(operator, file and line) as one of the \MutEarlierSampled{} of that run, though the destinations
may differ. \MutEarlierFailuresAbsent{} of the \MutEarlierFailures{} misses that were read are at no
location of the final sample.

Table~\ref{tab:repairs} lists every correction, one row each, in the order the chapters record them; the full record with commit hashes is in the accompanying repository. In the first column, \emph{figure} marks a correction made after an implausible or unwelcome figure was seen, \emph{panel} one from a red-team panel, \emph{review} one from a source review made with the assistant, \emph{harness} one that changed the mutation harness, and \emph{end-to-end check} a gap in the extractor that the suites of Section~\ref{sec:res-e2e} exposed, repaired afterwards with a test for each. The reviews and panels were made with the AI assistant that the declaration at the front of this thesis describes, and every defect they reported was reproduced against the code before it was repaired.

\begingroup\footnotesize\setlength{\tabcolsep}{2.5pt}
\begin{xltabular}{\textwidth}{@{}YYYY@{}}
\caption{The corrections made to the instruments.}
\label{tab:repairs}\\
\toprule
Found by & What was wrong & What changed & Effect on the figures \\
\midrule
\endfirsthead
\toprule
Found by & What was wrong & What changed & Effect on the figures \\
\midrule
\endhead
\bottomrule
\endfoot
\multicolumn{4}{@{}l}{\emph{Protocol and harness}} \\
One application's recall fell when the rule was tightened.
  & Test-source cross-check: a URL parameter in a test could match a literal route segment.
  & A parameter no longer matches a literal segment; a negative test is excluded.
  & The earlier recall cannot be regenerated; end-to-end recall there is now \EtoERateAlfresco{}. \\
\emph{Figure}: inspection of the flagged edges.
  & The first witness check showed a low precision for one application; all flagged edges were may edges from one over-approximated call site.
  & No extractor change. Over-approximated sites are reported apart.
  & The low figure described the measurement; the must tier had no failing edge. \\
\emph{Harness}: not stated beyond one application's misses.
  & A mutant whose target is absent from the baseline graph repeats a recovery gap counted elsewhere.
  & Such mutants form a separate category.
  & Most of that application's misses left the detection rate. \\
\emph{Figure}: a construct count that tripled.
  & The construct denominator counted the substring \texttt{routerLink}, which also matches \texttt{routerLinkActive}.
  & Constructs are counted by the independent route parser, so ratio and loss breakdown share one definition.
  & alfresco: the independent route parser finds \ConstructsAlfresco{}; the substring count was $\approx$\,three times larger. \\
\emph{Figure}: coverage reached 100 per cent (mifos).
  & The first independent route parser skipped wrapper calls, did not follow a lazily loaded module to its routing module, and dropped the prefix of lazy tables.
  & It reads wrapper calls, follows the routing module, and keeps the prefix.
  & mifos has \RoundTwoEarlierOracleMifosDeclared{} declared screens, \NodeAuditMatchedMifos{} with a node; the extractor was the same in both runs. \\
\emph{Figure}: a ratio that rose implausibly after attribution (mifos).
  & Must conversion was first must edges divided by constructs.
  & Computed per construct; edges whose screen set is a subtree are stored once (Section~\ref{sec:attribution}).
  & The rise came from how edges were stored, not from evidence. Definition: Section~\ref{sec:metrics}. \\
\emph{Harness}, \emph{review}, \emph{figure}: a review of the harness, and a detection rate that moved by more than twenty points.
  & The mutation harness patched the first textual match, judged detection on node and edge sets alone, and retargeted only the head of a commands array.
  & Patches at the construct's offset; detection also on the witness set; retargeting points at another declared route; guard mutants on a screen no edge enters are excluded; a rename counts when the renamed route appears.
  & realworld's detection rose sharply. Same seed, not the same sample: \MutEarlierSampled{} $\to$ \MutSampled{} mutants (earlier $\to$ final run), \MutOverlap{} at one location (operator, file, line; destinations may differ); \MutEarlierFailuresAbsent{}/\MutEarlierFailures{} earlier misses are not in the final sample. In the earlier run six paired disagreements favoured attribution; five were retarget mutants whose operator replaced only the head of the commands array, and with it corrected the baseline detects them too. \\
\emph{Harness}: reading the graph at the sites the rule let through.
  & The guard exclusion rule located a route by the nearest path literal before \texttt{canActivate}, which fails for children, empty paths, constants and template literals.
  & The route is found by brace-matching over source with comments and strings blanked; a route yielding no screen is excluded.
  & Four unobservable guard mutants had been scored as misses and explained as a defect of the graph; they are now excluded. The superseded run sampled \MutEarlierSampled{} mutants. \\
\emph{Harness}: a change in the extractor, not a figure.
  & Removed guards were misjudged: a child of a guarded parent lost the parent's guard, and a set comparison read three identical guard texts as unchanged.
  & A removal is detected when some edge loses a guard, excluded when no baseline edge carries it; guards are counted.
  & The sampled mutants are unchanged; how a sampled mutant is scored changes. \\
A table script stopped.
  & The construct denominator took constructs from commented-out links and code comments.
  & Comments are blanked before parsing, with the compiler's scanner for TypeScript.
  & One construct leaves the denominator in each of bitwarden and mifos; nothing else changes. \\
\addlinespace
\multicolumn{4}{@{}l}{\emph{The extractor and the deterministic tier}} \\
\emph{Panel}: red-team panels reading graph and source after commit \EarlierCommit{}; the witness check missed it.
  & A guard on a route a wrapper function returns was not read, so everything below was published as unconditional.
  & The guard set of a screen also reads routes a function returns.
  & At \EarlierCommit{}: \EarlierMustEdges{} must edges with \EarlierMustViolations{} witness failures, read from that commit's results. Reading the guard cut mifos to \RoundMustEdgesMifos{} on the intermediate version. \\
\emph{Panel}, as above.
  & A \texttt{canActivate} element that is a call, such as a guard factory, was skipped.
  & An element the reader cannot name gates the edge anyway.
  & bitwarden's deterministic edges into its root no longer read as unconditional. \\
Correcting the witness check itself.
  & It opened only the first of several constructs yielding one edge; later citations in bitwarden targeted \path{/vault}, which the extractor had not read.
  & It opens every cited site. An edge took the weakest modality of its constructs (replaced below).
  & Later citations are \WitnessMergedCitationShareAll{} of cited sites, \WitnessMergedCitationShareMax{} in the largest application. The intermediate tier held \RoundMustEdges{} must edges, \RoundMustClean{} passing, \RoundMustViolations{} failing. \\
\emph{Review}, then \emph{panel}: a review of every construct the intermediate tier cites (\MustCtxEdges{} edges, \MustCtxDistinct{} distinct constructs, cited \MustCtxConstructs{} times). The witness check cannot find this: it compares destinations, not whether a construct runs every time.
  & The demotion rules were incomplete: a link under \texttt{*ngIf}, in \texttt{@if}, a callback, a \texttt{catch}, or after \texttt{await} or an early return was must.
  & Demotion under every condition of Section~\ref{sec:model}; the panels added members no handler calls unconditionally, components under a template condition, ICU cases and addresses two routes can answer; a later review added code after a navigation that may start another.
  & \MustCtxConditional{}/\MustCtxConstructs{} citations ran under such a condition. \\
\emph{Review}, as above.
  & A guard on an empty-path child without a component was not read, and \texttt{canMatch}, \texttt{canActivateChild} and \texttt{canLoad} not at all.
  & Such a guard gates the screen above; the other three kinds are read.
  & \MustFateGuard{} of the intermediate tier's edges, all of bitwarden's, are guarded. \\
\emph{Panel}: a red-team panel.
  & The repair above gated an edge on every guard of its destination, demoting navigations the router lets through.
  & The extractor computes the guards the router runs for each source and destination: levels left, matched, entered, kept.
  & Navigations between screens under a shared guarded parent pass; mifos holds \CensusMustMifos{} must edges again. \\
\emph{Review}, as above.
  & A navigation to the address its source already displays was treated as occurring on every execution.
  & Such an edge is drawn per screen: must from the others, a may self-loop at the destination.
  & \MustFateSelfLoop{} of the intermediate tier's edges is such a self-loop. \\
Reading the edges that survived.
  & The weakest-modality rule discarded a certainty the edge had (realworld's Global Feed link runs whenever its page does).
  & An edge takes the modality of its most certain construct, its witness; other constructs record their own, and the witness check judges the must tier only on constructs that earn must.
  & The final tier is \CensusMust{} edges on \CensusMustSites{} constructs, a different set from the intermediate tier, not a smaller one. Superseded tables counted every cited construct as converted; current ones count a construct only where it earns must. \\
\addlinespace
\multicolumn{4}{@{}l}{\emph{The independent route parser and the denominators}} \\
The disagreements of the first node audit; each gap had scored the extractor against paths the independent route parser had not read.
  & It did not read a path written through a constant, an enum or \texttt{\$localize}, a first-party import through \texttt{baseUrl}, or a named route array.
  & It reads all three forms.
  & Before the last, it saw almost none of kubernetes/dashboard's screens and scored nearly every node the extractor minted there as unmatched. \\
The dead-link corrections.
  & It read Storybook fixtures as screens, a sibling application's tables via a tsconfig alias as this application's, and a path in an \texttt{Object.freeze} table as unresolvable.
  & It sets Storybook aside, ignores the sibling's tables, reads the frozen table. Dropping the frozen routes as unreadable was considered and rejected: they were readable one call away.
  & Declared sets moved in \RoundTwoSetMoved{}/\CensusApps{} applications; in \RoundTwoSetMovedCountSame{} of them the count stayed the same. \\
Fixing the provenance of the comparison.
  & A file from an interim revision of the parser, dropped into the table directory, was read as the earlier parser's output; the script could not detect this, and an absent file read as a denominator that did not move.
  & The earlier parser's own committed output is read from the repository.
  & The bitwarden decomposition built on that file is withdrawn. \\
Resolving the first node audit's disagreement.
  & The first-audit extractor reached route tables without their parent prefix and minted a node at \texttt{*} per catch-all.
  & The loss-list repairs: a table not reached, a path or destination named and not written (Section~\ref{sec:res-defects}).
  & Of \RoundTwoFirstRoundExtra{} nodes with no declared route, \RoundTwoFirstRoundTruncated{} were the tail of a declared route and \RoundTwoFirstRoundSynthetic{} catch-alls. The final version has \RoundTwoExtraNodes{} nodes matching no declared screen, \RoundTwoExtraNodesWildcard{} at \texttt{*}. \\
\emph{Not stated}; the check was made after the graph had been consulted.
  & The grounds that exclude a construct from the denominator were tested after the graph had been consulted.
  & The grounds are decided from the construct's own text first; the comparison refuses to print if a construct changes ground between runs.
  & The after column had been over more sites than the before column; both are now over the same \RoundTwoInScope{} sites. \\
\addlinespace
\multicolumn{4}{@{}l}{\emph{Findings of the runtime route oracle}} \\
The oracle returned two alfresco addresses the framework answers with its catch-all; both static readers read the subtree alike, so 100 per cent coverage hid it.
  & A route with a named outlet could take the primary address from a sibling, and routes beneath it were minted at the primary address.
  & Outlet context is carried down the walk; no screen is minted beneath a named outlet.
  & The extractor declares \ProbeInvented{} screens at addresses the framework mounts nothing at. What remains limits recall. \\
The oracle against the route set as the static readers held it; both readers had arrived at the defect independently.
  & A lazy route naming a module through a barrel was followed one hop only, and the table scanned at the root.
  & Both follow the chain to the table, bounded, refusing at each hop a file the route sits in, a file another route loads, and a file whose tables are registered at the root.
  & \ProbeBeforeNotMounted{} addresses in \ProbeBeforeWhere{} were wrong, the two alfresco ones among them. Node coverage had been 100 per cent against a denominator with the same wrong addresses. \\
Instrumenting the predicate (commit \texttt{d40c8dc}).
  & The test keeping a lazy route from adopting the application's root table read \texttt{RouterModule.\allowbreak forRoot}, so it fired nowhere in the corpus.
  & Both tools mark root tables per array, from the argument of either registration API.
  & A guard that never fires had been credited with a property nothing showed it had. \\
\addlinespace
\multicolumn{4}{@{}l}{\emph{The dead-link census}} \\
\emph{Refutation attempts} on the dead-link finding.
  & An earlier correction followed only the redirecting kind of catch-all route.
  & Both kinds are followed, the edge is held at may, and where the catch-all sent the URL is recorded.
  & The correction: \checkmark{} for the application it was written against, \xmark{} for the two needing the other kind. \\
\emph{Refutation attempts}.
  & Resolving a target and calling a link dead were conflated; bitwarden's authentication paths went through an unread enum, so the wildcard forwarded them all.
  & The enum and wrapper call are read; each application notes when its route set is a lower bound; each caveated candidate is put to the framework's route configuration.
  & The wildcard had raised that application's recovered figure. Withdrawn: \RoundTwoFrameworkMountsIt{} (the framework serves a screen), \RoundTwoFrameworkPrefix{} (the matching route places none), \RoundTwoFrameworkRedirects{} (a redirect answers); \RoundTwoWildcardCaveated{} keep the caveat. \\
\emph{Refutation attempts}.
  & Unreadable routes were taken as undeclared: a spread-built array, ghostfolio's \texttt{matcher}, alfresco's tsconfig \texttt{paths} alias.
  & Both tools record the holes they cannot read; no undeclared claim is counted inside one.
  & ghostfolio's \path{/admin/users} had been reported as a confirmed dead link, and two live alfresco links as dead; all three constructs are read now. \\
\emph{Refutation attempts}.
  & The route-set test read only the first quoted segment of a commands array, testing a prefix as the target.
  & The whole array is read as a pattern, computed segments as holes, matched pattern to pattern.
  & It cleared alfresco's two candidates and found further sites in mifos; the census holds \RoundTwoDeadLinkCandidates{} candidate sites. \\
\addlinespace
\multicolumn{4}{@{}l}{\emph{Mutation sweep and end-to-end check}} \\
\emph{Mutation sweep} of the demotion rules, at commit \SweepEarlierCommit{}.
  & On its first run in the final revision the sweep found rules the router fixes had added without a test.
  & Tests were written for most survivors; the rest were retired as equivalent, since no input tells them from their absence.
  & \SweepEarlierSurvived{}/\SweepEarlierMutants{} mutants survived; \SweepEarlierCovered{} now have tests, \SweepEarlierRetired{} were retired. At \SweepCommit{} every one of \SweepMutants{} deletions fails some test. The tests exposed a relative link that climbs above the root. \\
\emph{Panel}: a review panel checking the scorer against the preregistration.
  & The end-to-end scorer compared a redirected navigation on the landed address, against the preregistration.
  & The scorer applies the preregistered rule.
  & mifos then scored \EtoEMifosHitUnfixed{}/\EtoEMifosPairs{} pairs, \EtoEMifosRedirect{} of them redirected; the landed-target figure had been \EtoEMifosHitAfterRedirectsUnfixed{}/\EtoEMifosPairsAfterRedirects{}. \\
\emph{End-to-end check}: mifos's suite missed transitions to action screens.
  & A template literal inside a commands array (\texttt{navigate([\textasciigrave{}actions/}\allowbreak\texttt{\$\{name\}\textasciigrave{}])}) was read as a variable, so its target was recorded as unresolved.
  & Its static text is read, each segment holding an interpolation a parameter.
  & The client, group and savings views reach their action screens. \\
\emph{Figure}: after that repair the re-measured mifos graph joined every pair of screens.
  & A target prefix was cut at its first parameter, so the notifications tray's prefix, which opens with two parameters, was empty and the shell site reached every route.
  & Only trailing parameters are dropped from a target prefix; a parameter inside it matches any one segment.
  & The tray reaches the account screens its targets name; possible pairs the mifos graph holds, before $\to$ after the repairs: \EtoEMifosGraphPairsUnfixed{} $\to$ \EtoEMifosGraphPairs{}/\EtoEMifosPossiblePairs{}. \\
\emph{End-to-end check}: the same.
  & \texttt{relativeTo} passed in an options variable was not evaluated, so the loan view's actions were resolved from the root.
  & A \texttt{const} initialised with an object literal is read as that literal.
  & The loan view reaches its action screens. \\
\emph{End-to-end check}: mifos's suite missed tab transitions.
  & A relative target was resolved against the screen it fires on; a link in a parent's template, rendered with each child tab, named nothing mounted.
  & It is resolved against the route that renders the component holding it (Section~\ref{sec:anchored}); the witness check applies the same rule (harness \HarnessVersion{}).
  & With the two rows above, mifos suite pairs held: \EtoEMifosHitUnfixed{} $\to$ \EtoEMifosHit{}/\EtoEMifosPairs{}; every figure that depends on the extractor was re-measured. \\
\end{xltabular}
\par\smallskip{\footnotesize\noindent\emph{Note.} \emph{Found by} names the method as the main text does; \emph{figure}, \emph{panel}, \emph{review}, \emph{harness} and \emph{end-to-end check} are explained in the text above.}
\endgroup

Several corrections in the protocol and harness group were found by a figure that looked implausibly good or bad (marked \emph{figure}); each proved a property of the measurement and not of the extractor.

\subsection{Reproducibility}
\label{sec:reproducibility}

Every figure is generated by scripts at one commit of the extractor and written into the thesis as a
macro. \texttt{REPRODUCE.md} in the repository lists the scripts and the commit each table comes from. The version evaluated throughout is commit \texttt{f26a302}, whose extractor last changed at \texttt{2aac284}, with harness 1.10.0 and mutation harness 1.9.1; the demotion-rule sweep and the may-reason census were recorded at the same commit. The end-to-end figure before the repairs of Table~\ref{tab:repairs} comes from the graph of the previously evaluated commit, \texttt{c42f220}, kept beside the results. Two groups are historical and were not re-measured: the rendered-interface check (Table~\ref{tab:scope}), and the runtime route oracle, run on 29 and 30 October against a graph whose extractor commit was not recorded and whose node total is its only link to the final graph.

\section{Summary}

The corpus is six applications at pinned commits, selected by explicit criteria that were not
applied uniformly, and the extractor was developed on it. Every ratio has one name and one definition in
Table~\ref{tab:ratios}. The five methods are each blind where the others see; the reproducibility
arrangements cover the first four fully and the fifth, the runtime route oracle, only in part. The
corrections that led to the final protocol are listed in Table~\ref{tab:repairs}.
